%%%%%%%%%%%%%%%%%%%%%%%%%%%%%%%%%%%%%%%%%%%%%%%%%%%%%%%%%%%%%%%%%%%%%%%%%%%%%%%%%%
% Career in MLOps (Sep 2026): what the work is, which roles exist, which skills to build, how to start and prepare.
% Session 19 only. Certification and course names change: check the vendor pages before teaching.
\begin{frame}[fragile]\frametitle{}
\begin{center}
{\Large Building a Career in MLOps}
\end{center}
\end{frame}

%%%%%%%%%%%%%%%%%%%%%%%%%%%%%%%%%%%%%%%%%%%%%%%%%%%
\begin{frame}[fragile]\frametitle{Why MLOps Is a Career}
\begin{itemize}
\item Many models stall between the notebook and the real world: the Netflix story
\item Someone has to build the road to production
\begin{itemize}
\item Check the data, package the model, deploy it, monitor it, retrain it
\end{itemize}
\item MLOps combines three skill sets: machine learning, software engineering and operations
\item The same skills carry over to newer systems, such as applications built on language models
\end{itemize}
\end{frame}

%%%%%%%%%%%%%%%%%%%%%%%%%%%%%%%%%%%%%%%%%%%%%%%%%%%
\begin{frame}[fragile]\frametitle{Roles Around MLOps}
\begin{center}
\small\renewcommand{\arraystretch}{1.35}
\begin{tabular}{>{\raggedright\arraybackslash}p{0.3\linewidth}>{\raggedright\arraybackslash}p{0.62\linewidth}}
\hline
\textbf{Role} & \textbf{What you do} \\
\hline
MLOps engineer & build and run the pipelines for deployment, monitoring and retraining \\
Machine learning engineer & train models and make them ready for production \\
ML platform engineer & build the shared tools that many ML teams use \\
Data engineer & build the data pipelines that feed the models \\
DevOps or reliability engineer & run the infrastructure; add ML knowledge \\
\hline
\end{tabular}
\end{center}
Job titles differ between companies: read the description, not the title.
\end{frame}

%%%%%%%%%%%%%%%%%%%%%%%%%%%%%%%%%%%%%%%%%%%%%%%%%%%
\begin{frame}[fragile]\frametitle{Where the Opportunities Are}
\begin{itemize}
\item Employers
\begin{itemize}
\item Product companies and start-ups that run ML-based services
\item IT services and consulting firms, and global capability centres of large companies
\item Manufacturing, automotive and energy companies
\end{itemize}
\item Work for a mechanical engineer: predictive maintenance, quality inspection, simulation surrogate models, energy systems
\item Your domain knowledge is an advantage: you know the sensors, the failure modes and the physics
\item A common way in: start as a data, ML or DevOps engineer, then grow towards MLOps
\end{itemize}
\end{frame}

%%%%%%%%%%%%%%%%%%%%%%%%%%%%%%%%%%%%%%%%%%%%%%%%%%%
\begin{frame}[fragile]\frametitle{Skills You Need: Tools}
\begin{center}
\small\renewcommand{\arraystretch}{1.35}
\begin{tabular}{>{\raggedright\arraybackslash}p{0.27\linewidth}>{\raggedright\arraybackslash}p{0.6\linewidth}}
\hline
\textbf{Area} & \textbf{What to learn} \\
\hline
Programming & Python, testing, building an API (FastAPI, as in this session) \\
Linux and git & the shell, scripts, branches, pull requests \\
Machine learning & the basics from this course: data, models, evaluation \\
Data & SQL, data validation, data pipelines \\
Containers and cloud & Docker, one cloud provider, later Kubernetes \\
Automation & CI/CD pipelines, infrastructure as code (for example Terraform) \\
Monitoring & logs, metrics and drift checks (for example Prometheus, Grafana, Evidently) \\
\hline
\end{tabular}
\end{center}
\end{frame}

%%%%%%%%%%%%%%%%%%%%%%%%%%%%%%%%%%%%%%%%%%%%%%%%%%%
\begin{frame}[fragile]\frametitle{Skills You Need: Beyond Tools}
\begin{itemize}
\item Ownership: someone is called when the service breaks, and it is you
\item Communication: you connect data scientists, software developers and the business
\item Cost awareness: compute, storage and licenses all cost money
\item Security and privacy: know who may see the data and the model
\item Debugging: read logs and documentation patiently
\item Domain knowledge: understand the machine or the process behind the data
\end{itemize}
\end{frame}

%%%%%%%%%%%%%%%%%%%%%%%%%%%%%%%%%%%%%%%%%%%%%%%%%%%
\begin{frame}[fragile]\frametitle{A Possible Roadmap}
A starting guide, not a rule: the pace depends on the time you have.
\begin{enumerate}
\item Months 1 to 3: Python, Linux, git, and the machine learning basics from this course
\item Months 3 to 6: Docker, REST APIs, SQL, and one cloud provider
\item Months 6 to 9: CI/CD, experiment tracking and a model registry (for example MLflow), pipeline orchestration
\item Months 9 to 12: monitoring, Kubernetes, and a capstone project that ties everything together
\end{enumerate}
\end{frame}

%%%%%%%%%%%%%%%%%%%%%%%%%%%%%%%%%%%%%%%%%%%%%%%%%%%
\begin{frame}[fragile]\frametitle{How to Prepare: Show Your Work}
\begin{itemize}
\item Build one public project from end to end: train, serve, containerize, test and monitor a model
\begin{itemize}
\item This session's walkthrough is a good start
\end{itemize}
\item Write a clear README and draw the architecture on one page
\item Finish a free course with a project, for example the MLOps Zoomcamp
\item Fix a bug or improve the documentation of an open-source MLOps tool
\item Apply for internships and join hackathons
\item Practise explaining your design choices: why this tool, what would break, what it costs
\end{itemize}
\end{frame}

%%%%%%%%%%%%%%%%%%%%%%%%%%%%%%%%%%%%%%%%%%%%%%%%%%%
\begin{frame}[fragile]\frametitle{Certificates and Courses}
\begin{itemize}
\item Cloud certificates
\begin{itemize}
\item AWS Certified Machine Learning Engineer, Associate
\item Google Cloud Professional Machine Learning Engineer
\item Microsoft: the Azure Data Scientist exam DP-100 was retired in June 2026; check its successor (AI-300)
\end{itemize}
\item Kubernetes: the Certified Kubernetes Administrator (CKA)
\item Courses and books
\begin{itemize}
\item MLOps Zoomcamp (DataTalks.Club, free)
\item Designing Machine Learning Systems (Chip Huyen)
\item Machine Learning in Production (DeepLearning.AI)
\end{itemize}
\item A certificate helps you get an interview; a project proves you can do the work
\end{itemize}

{\tiny Exams and courses change: check the provider's page before you commit.}
\end{frame}
